\section{Memory Management in Go}

\subsection{Introduction}

Memory management is a fundamental aspect of programming language design. It concerns the allocation, use, organisation and reclamation of memory required to store program data during execution. Effective memory management is important for application performance, reliability, scalability and resource utilisation.\\

The Go programming language provides automatic memory management through a garbage collector, together with language features and runtime mechanisms that manage memory allocation and object lifetimes. Unlike languages such as C and C++, Go does not ordinarily require programmers to explicitly allocate and free memory using operations equivalent to \texttt{malloc} and \texttt{free}. Instead, the Go compiler and runtime cooperate to determine where values can be stored and when allocated objects are no longer required.\\

Go's memory management system incorporates several mechanisms:

\begin{itemize}
    \item Automatic allocation of variables and dynamically sized data structures.
    \item Compiler escape analysis to determine whether values can remain on the stack or must be allocated on the heap.
    \item A concurrent, tracing garbage collector that identifies and reclaims unreachable heap objects.
    \item Runtime management of goroutine stacks and memory allocation.
    \item Explicit control over object references through pointers, slices, maps and interfaces.
    \item Allocation and reuse of memory through runtime allocators and size classes.
\end{itemize}

Although Go automates much of the memory-management process, developers remain responsible for designing data structures, controlling object lifetimes, avoiding unnecessary allocations and understanding how references affect garbage collection.

\subsection{The Go Memory Model}

A useful starting point for understanding memory management in Go is to distinguish between the stack and the heap. These terms describe different storage regions and allocation strategies, although the precise placement of a particular value is determined by compiler and runtime implementation details.

\subsubsection{Stack Memory}

The stack is used for function execution state, including local variables, function arguments, return-related information and temporary values. In Go, each goroutine has its own stack.\\

Goroutine stacks begin relatively small and can grow when more stack space is required. The runtime manages this growth automatically. Consequently, developers do not ordinarily specify a fixed stack size for each goroutine.

Consider the following example:

\begin{lstlisting}
package main

import "fmt"

func calculate() int {
    x := 10
    y := 20
    return x + y
}

func main() {
    result := calculate()
    fmt.Println(result)
}
\end{lstlisting}

The variables \texttt{x} and \texttt{y} are local to the \texttt{calculate} function. If the compiler determines that their values do not need to outlive the function call and that stack allocation is appropriate, it may keep them in stack storage or optimise them into registers.\\

When the function returns, its stack frame is no longer active. However, it is important to understand that not every local variable is necessarily stored on the stack. The compiler may eliminate a variable entirely, keep its value in a register, or allocate it on the heap if its lifetime or usage requires that placement.\\

Stack allocation is generally inexpensive because it can often be implemented through adjustments to the stack pointer rather than a general heap-allocation operation.

\subsubsection{Heap Memory}

The heap is used for objects whose lifetimes or storage requirements make stack allocation unsuitable. Heap objects are managed by the Go runtime and are subject to garbage collection when they become unreachable.\\

For example:

\begin{lstlisting}
package main

import "fmt"

type User struct {
    Name string
    Age  int
}

func createUser() *User {
    user := &User{
        Name: "Alice",
        Age:  30,
    }

    return user
}

func main() {
    user := createUser()
    fmt.Println(user.Name)
}
\end{lstlisting}

The function returns a pointer to a \texttt{User} value. The value must remain valid after \texttt{createUser} returns. The compiler's escape analysis will ordinarily determine that the object must outlive the function's stack frame and allocate it on the heap.\\

The key principle is that Go allows a pointer to refer safely to a value beyond the function in which that value was created. The compiler and runtime arrange the storage accordingly.\\

Heap allocation generally has more overhead than a simple stack allocation. Heap objects also contribute to garbage-collection work, so excessive allocation can increase memory consumption and CPU usage.\\

However, heap allocation is not inherently undesirable. It is necessary for many useful programming patterns, including returning references to data, sharing objects between goroutines and maintaining dynamically sized data structures.

\subsection{Automatic Memory Allocation}

Go allocates storage for values as required by the program. Allocation may occur implicitly when declaring variables, creating composite literals, constructing slices or maps, or using operations such as \texttt{new} and \texttt{make}.

\subsubsection{Variable Declarations}

A variable declaration establishes a typed variable:

\begin{lstlisting}
var count int
var name string = "Alice"
active := true
\end{lstlisting}

The variables are initialised according to their types and assigned values. The compiler determines their storage and may optimise their representation.\\

For types such as integers and Booleans, the value can often be represented directly in a machine word or register. More complex types may contain references to other storage.\\

For example, a string is conceptually a descriptor containing a pointer to its character data and a length. A slice contains a pointer to an underlying array, a length and a capacity. A map is a reference to a runtime-managed data structure.\\

These representations help explain why copying a slice or map value does not necessarily copy all the data associated with it.

\subsubsection{The \texttt{new} Function}

The built-in \texttt{new} function allocates storage for a value of a specified type, initialises it to that type's zero value and returns a pointer to it.

\begin{lstlisting}
package main

import "fmt"

func main() {
    p := new(int)

    fmt.Println(*p) // 0

    *p = 42
    fmt.Println(*p) // 42
}
\end{lstlisting}

The expression \texttt{new(int)} returns a value of type \texttt{*int}. The pointed-to integer is initially zero.\\

The language specification does not guarantee that every value created using \texttt{new} is allocated on the heap. Compiler optimisation and escape analysis determine the appropriate storage location.\\

The \texttt{new} function is therefore not a direct equivalent of an unconditional heap-allocation operation in C or C++.\\

\subsubsection{The \texttt{make} Function}

The built-in \texttt{make} function initialises values of certain built-in reference-oriented types, namely slices, maps and channels.\\

For example:

\begin{lstlisting}
numbers := make([]int, 5, 10)
lookup := make(map[string]int)
messages := make(chan string, 4)
\end{lstlisting}

The slice has an initial length of five and a capacity of ten. The map is initialised and ready to store key-value pairs, while the channel is created with a buffer capacity of four values.\\

Unlike \texttt{new}, which returns a pointer to a zero-valued object, \texttt{make} returns an initialised slice, map or channel value.\\

The internal storage required by these values is managed by the runtime. For example, a slice may require an underlying array, while a map requires runtime-managed storage for its entries.\\

\subsection{Escape Analysis}

Escape analysis is a compiler technique used to determine whether a value can safely remain in a limited lifetime region, such as a function's stack frame, or must be placed in storage that outlives that region.\\

Go performs escape analysis during compilation. The resulting decisions influence allocation behaviour and garbage-collection overhead.\\

\subsubsection{A Value That Does Not Escape}

Consider the following function:

\begin{lstlisting}
func add(a, b int) int {
    result := a + b
    return result
}
\end{lstlisting}

The function returns an integer value rather than a pointer to its local variable. The compiler can generally keep the relevant values in registers or stack storage without requiring a heap allocation.\\

In practice, an optimiser may remove the variable entirely and calculate the result directly.

\subsubsection{A Value That Escapes}

Now consider:

\begin{lstlisting}
type Record struct {
    ID int
}

func createRecord() *Record {
    record := Record{ID: 100}
    return &record
}
\end{lstlisting}

The function returns the address of a local variable. Since the caller may retain the returned pointer after the function has returned, the value must remain valid beyond the lifetime of the function's stack frame.\\

The compiler can therefore arrange for \texttt{record} to be allocated on the heap.\\

This demonstrates an important distinction: returning the address of a local variable is safe in Go. The compiler does not leave the returned pointer referring to invalid stack memory.\\

\subsubsection{Escape Analysis Is Not Simply About Returning Pointers}

A value can escape even when its address is not returned directly. For example, storing a pointer in a global variable or placing it in a data structure whose lifetime exceeds the current function may require heap allocation.\\

Conversely, taking the address of a variable does not automatically force a heap allocation. If the compiler can prove that the pointer does not outlive the appropriate storage region, the value may remain on the stack.\\

The compiler can also make different allocation decisions depending on inlining, optimisation settings, Go versions and the surrounding program.\\

Escape-analysis decisions can be inspected using compiler diagnostic options:

\begin{lstlisting}
go build -gcflags="-m" .
\end{lstlisting}

For more detailed diagnostics, the following form can be used:

\begin{lstlisting}
go build -gcflags="-m=2" .
\end{lstlisting}

These options ask the compiler to report optimisation and escape-analysis information. The output can help identify values that escape to the heap and functions that are inlined.\\

Such diagnostics are implementation-specific and may change between compiler versions. They should be used to investigate performance rather than as a substitute for benchmarking.\\

\subsection{Garbage Collection}

Garbage collection is the automatic reclamation of memory occupied by objects that are no longer reachable by the program.\\

Go's garbage collector is a tracing, concurrent, non-generational mark-and-sweep collector. Its principal task is to identify heap objects that remain reachable from a set of root references and reclaim objects that cannot be reached.\\

The collector is integrated into the Go runtime and is designed to perform most of its work concurrently with application goroutines. Certain phases require brief stop-the-world coordination, during which application goroutines are paused.

\subsubsection{Reachability}

An object is reachable if the program can access it through a chain of references beginning at a garbage-collection root.\\

Important roots include relevant global variables, active goroutine stacks and runtime-maintained references. The exact root set and implementation details depend on the runtime.\\

Consider the following example:

\begin{lstlisting}
type Node struct {
    Value int
    Next  *Node
}

func main() {
    first := &Node{Value: 1}
    second := &Node{Value: 2}

    first.Next = second

    fmt.Println(first.Value)
}
\end{lstlisting}

The variable \texttt{first} refers to a node, and that node refers to \texttt{second}. As long as the relevant reference chain remains reachable, both nodes must be preserved.\\

If the program removes all reachable references to the two nodes, the garbage collector may eventually reclaim their storage.\\

The collector does not normally reclaim an object merely because it has not been used recently. It determines reachability, not semantic usefulness. An object that is no longer useful to the application may remain reachable through an unintended reference and therefore cannot be collected.\\

\subsubsection{Mark-and-Sweep Collection}

Go's garbage collector broadly operates using marking and sweeping.

\begin{enumerate}
    \item \textbf{Marking:} The collector identifies reachable objects by traversing references from the root set.
    \item \textbf{Sweeping:} Memory associated with unreachable objects is made available for reuse by the runtime.
\end{enumerate}

During concurrent marking, the runtime uses write barriers to maintain the correctness of the marking process while application goroutines continue modifying pointers.\\

A write barrier is a runtime mechanism associated with pointer writes. It helps ensure that objects reachable during concurrent marking are not incorrectly treated as unreachable because the object graph changes while the collector is operating.\\

The collector's concurrent design reduces the time applications spend paused for garbage collection, although collection still consumes CPU resources and can introduce latency.\\

The garbage collector is not a reference-counting system. It does not ordinarily reclaim an object immediately when the number of references to that object becomes zero. Instead, objects are reclaimed as part of a collection cycle when the tracing process determines that they are unreachable.\\

\subsubsection{An Object Becoming Eligible for Collection}

Consider:

\begin{lstlisting}
func process() {
    data := make([]byte, 1024*1024)

    data[0] = 42

    fmt.Println(data[0])
}
\end{lstlisting}

The slice refers to an underlying array containing approximately one mebibyte of byte storage. When the underlying storage is no longer reachable, it becomes eligible for garbage collection.\\

Returning from \texttt{process} does not guarantee that the memory will immediately be reclaimed. The runtime may retain the storage for reuse, and the next collection may occur later.\\

This distinction is important because the amount of memory allocated by a program, the amount of live data and the amount of memory retained by the runtime are different quantities.\\

\subsection{Concurrent Garbage Collection and the Runtime}

Go's garbage collector is designed to limit interference with application execution. Much of the marking work takes place concurrently with the execution of goroutines.\\

The runtime coordinates collection through a combination of scheduling, marking, write barriers and brief stop-the-world phases. A collection cycle typically includes preparation, concurrent marking, mark termination, sweeping and associated background work.\\

The exact sequence and implementation details can vary between Go releases, but the principal design objective is to reclaim unreachable memory without requiring the application to stop for the entire duration of the collection.\\

Concurrent collection has several consequences:

\begin{itemize}
    \item Application goroutines can continue executing during much of the marking phase.
    \item The collector consumes CPU time that would otherwise be available to application code.
    \item Pointer-rich data structures generally require more tracing work than data structures containing large amounts of pointer-free data.
    \item The amount of live heap data influences the cost of collection.
    \item Allocation rate affects how frequently collection is needed and how much memory is allocated between cycles.
\end{itemize}

A program that creates large numbers of short-lived objects can experience significant garbage-collection overhead even if its long-term live-data size is relatively small.\\

Reducing unnecessary allocations can therefore improve both throughput and latency. Nevertheless, optimisation should be based on profiling because changes that reduce allocation counts may increase code complexity or create other performance costs.

\subsection{Garbage-Collection Tuning}

Go provides mechanisms for influencing garbage-collection behaviour. One important control is the \texttt{GOGC} environment variable or its corresponding runtime setting.\\

The default value of \texttt{GOGC} is normally 100. Conceptually, this controls the heap-growth target relative to the amount of live heap data at the end of the previous collection, subject to runtime accounting and minimum-heap rules.\\

A lower value generally triggers collection more frequently, potentially reducing memory growth while increasing garbage-collection CPU overhead. A higher value generally permits more heap growth between collections, potentially reducing collection frequency at the cost of higher memory usage.\\

For example, a program can set the target programmatically:

\begin{lstlisting}
package main

import "runtime/debug"

func main() {
    debug.SetGCPercent(200)

    // Application logic follows.
}
\end{lstlisting}

The value 200 requests a higher heap-growth target than the default setting. It does not mean that the program will use exactly twice as much memory, nor does it establish a strict upper bound on memory consumption.\\

The garbage collector can also be disabled through a negative GC percentage. This should be used cautiously because allocations may continue to accumulate and cause substantial memory consumption.\\

For applications with explicit memory limits, Go also provides a soft memory limit through \texttt{GOMEMLIMIT} or \texttt{runtime/debug.SetMemoryLimit}. The limit influences garbage-collection pacing and the runtime's memory-management decisions. It is not a hard guarantee that total process memory can never exceed the specified value, particularly when the application requires more memory to make progress or when memory outside the managed heap contributes to the process footprint.\\

Memory tuning should be guided by measurements of allocation rate, live heap size, garbage-collection CPU usage, latency and overall process memory.

\subsection{Pointers and Memory Safety}

Pointers provide a way to refer to values stored elsewhere in memory. Go supports pointers but deliberately excludes general pointer arithmetic from ordinary safe Go code.\\

For example:

\begin{lstlisting}
package main

import "fmt"

func main() {
    value := 42
    pointer := &value

    fmt.Println(*pointer)

    *pointer = 99
    fmt.Println(value)
}
\end{lstlisting}

The expression \texttt{\&value} obtains the address of the variable, while \texttt{*pointer} dereferences the pointer to read or modify the pointed-to value.\\

Go's type system and runtime eliminate many common memory errors associated with manual memory management. In ordinary safe Go code, a pointer does not become invalid merely because the function that created the pointed-to value has returned.\\

However, Go is not immune to every memory-related bug. Incorrect pointer use can still cause a panic through a nil dereference, and unsafe operations can violate the language's memory-safety guarantees. Data races, excessive memory retention and resource leaks are also possible.\\

The \texttt{unsafe} package permits certain low-level operations that bypass parts of Go's usual type-safety protections. It should be used only when necessary and with a clear understanding of the associated risks.\\

\subsection{Slices and Their Memory Behaviour}

A slice is a descriptor referring to a contiguous segment of an underlying array. It contains a pointer to the array, a length and a capacity.\\

The length specifies how many elements are accessible through the slice, while the capacity specifies how many elements can be accessed through reslicing before a new underlying array is required.\\

Consider:

\begin{lstlisting}
package main

import "fmt"

func main() {
    values := make([]int, 3, 5)

    values[0] = 10
    values[1] = 20
    values[2] = 30

    fmt.Println(len(values)) // 3
    fmt.Println(cap(values)) // 5

    values = append(values, 40)

    fmt.Println(values)
}
\end{lstlisting}

The slice initially has a length of three and a capacity of five. Appending one element can use the existing underlying array because spare capacity is available.\\

If an append operation exceeds the current capacity, Go allocates a new underlying array, copies the existing elements and returns a slice referring to the new storage. The precise capacity growth strategy is an implementation detail and should not be relied upon as a fixed formula.\\

A significant memory-management issue arises when a small slice retains a reference to a much larger underlying array:

\begin{lstlisting}
func extractPrefix(data []byte) []byte {
    return data[:10]
}
\end{lstlisting}

If the returned slice remains reachable, its reference to the underlying array can keep the entire array alive, even though only ten bytes are accessible through the slice.\\

If retaining the large array is undesirable, the required portion can be copied into a new slice:

\begin{lstlisting}
func extractPrefix(data []byte) []byte {
    result := make([]byte, 10)
    copy(result, data[:10])
    return result
}
\end{lstlisting}

This approach allows the original array to become eligible for collection once no other references remain. It also introduces a new allocation and a copy operation, so the trade-off should be evaluated according to the application's workload.

\subsection{Maps and Memory Allocation}

A Go map is a built-in data structure that associates keys with values. Its internal representation is managed by the runtime.\\

For example:

\begin{lstlisting}
users := make(map[int]string)

users[1] = "Alice"
users[2] = "Bob"

fmt.Println(users[1])
\end{lstlisting}

The \texttt{make} function creates an initialised map that is ready for use. The runtime manages the storage required for its entries and may allocate additional storage as the map grows.\\

When the approximate number of entries is known, a capacity hint can be supplied:

\begin{lstlisting}
users := make(map[int]string, 1000)
\end{lstlisting}

The hint can reduce the need for repeated growth-related allocations in some workloads. It is not a guarantee that exactly 1000 entries will be stored without further allocation, nor does it establish a maximum map size.\\

Removing entries from a map does not guarantee that the runtime immediately returns all associated memory to the operating system. Internal storage may remain available for reuse.\\

A map should be allowed to become unreachable when it is no longer required. If a long-lived map retains references to large objects, those objects may remain reachable even when they are no longer useful to the application.\\

\subsection{Memory Management in Concurrent Programs}

Go's goroutines enable concurrent execution, and the runtime manages their stacks independently. Each goroutine begins with a relatively small stack that can grow when necessary.\\

The runtime may relocate a goroutine's stack when growing or shrinking it. Consequently, Go code should not assume that a goroutine's stack has a fixed size or that the address of a stack variable remains unchanged throughout the goroutine's lifetime.\\

The garbage collector accounts for references on goroutine stacks when determining which heap objects remain reachable.\\

Consider:

\begin{lstlisting}
package main

import (
    "fmt"
    "sync"
)

func main() {
    var wg sync.WaitGroup

    for i := 0; i < 5; i++ {
        wg.Add(1)

        go func(value int) {
            defer wg.Done()
            fmt.Println(value)
        }(i)
    }

    wg.Wait()
}
\end{lstlisting}

Each goroutine receives its own argument value. The \texttt{sync.WaitGroup} ensures that the main goroutine waits for all five goroutines to finish.\\

The lifetime of data used by a goroutine must be considered carefully. If a goroutine retains a pointer to an object, the object may remain reachable for as long as that reference remains live.\\

Furthermore, a goroutine that blocks indefinitely can retain stack variables and referenced heap objects. Such a goroutine may therefore cause memory to remain allocated even if the application no longer needs the work it was performing.\\

Garbage collection does not terminate goroutines or release resources held by goroutines that never finish. Applications should design explicit cancellation and shutdown mechanisms, commonly using channels and contexts, to ensure that concurrent work can terminate when it is no longer required.\\

Concurrent access to shared memory also requires appropriate synchronisation. Garbage collection does not prevent data races. Mutexes, channels, atomic operations and other synchronisation mechanisms are used to coordinate access to shared mutable state.\\

\subsection{Memory Retention and Resource Leaks}

Automatic garbage collection removes the need for explicit deallocation in ordinary Go programs, but it does not eliminate memory leaks in the broader sense.\\

A memory leak can occur when a program retains references to objects that it no longer needs. Because the objects remain reachable, the garbage collector cannot reclaim them.\\

For example:

\begin{lstlisting}
var cache = make(map[string][]byte)

func store(key string, data []byte) {
    cache[key] = data
}
\end{lstlisting}

This global cache retains every stored entry until an entry is removed or the map becomes unreachable. If the application continuously adds entries without a retention policy, memory consumption can grow indefinitely.\\

A bounded cache, explicit deletion, expiration policies, or another storage strategy may be necessary to control growth.\\

Resources such as open files, sockets, database connections, and operating-system handles require separate consideration. These are not equivalent to ordinary heap objects and should be released explicitly when no longer required.\\

For example:

\begin{lstlisting}
file, err := os.Open("data.txt")
if err != nil {
    return err
}
defer file.Close()
\end{lstlisting}

The deferred call closes the file when the surrounding function returns. The garbage collector should not be relied upon as the primary mechanism for closing files or releasing other external resources.

\subsection{Reducing Unnecessary Allocations}

Reducing unnecessary allocations can improve throughput, lower garbage-collection pressure, and reduce memory usage. However, allocation optimisation should preserve correctness and maintainability.\\

Several common techniques are available.

\subsubsection{Preallocating Slices}

If the approximate number of elements is known, a slice can be created with sufficient capacity:

\begin{lstlisting}
items := make([]string, 0, 1000)

for i := 0; i < 1000; i++ {
    items = append(items, fmt.Sprint(i))
}
\end{lstlisting}

Preallocating capacity may reduce the number of underlying-array growth operations and associated copies.\\

The capacity should be chosen sensibly. Excessively large preallocations can waste memory, especially when many instances are created concurrently.

\subsubsection{Reusing Buffers}

The \texttt{bytes.Buffer} type provides reusable byte storage:

\begin{lstlisting}
var buffer bytes.Buffer

buffer.WriteString("Hello, ")
buffer.WriteString("world!")

fmt.Println(buffer.String())

buffer.Reset()
\end{lstlisting}

Calling \texttt{Reset} sets the buffer's logical contents to empty while retaining its underlying storage for reuse. This can reduce allocations when a buffer is repeatedly used for similar tasks.\\

However, a buffer that once held a very large value may retain a correspondingly large backing array. Reusing such buffers indefinitely can increase memory retention, so applications may need to discard oversized buffers.\\

\subsubsection{Using \texttt{sync.Pool}}

The \texttt{sync.Pool} type provides a mechanism for temporarily reusing objects between independent operations:

\begin{lstlisting}
var bufferPool = sync.Pool{
    New: func() any {
        return new(bytes.Buffer)
    },
}

func process() {
    buffer := bufferPool.Get().(*bytes.Buffer)
    buffer.Reset()

    defer bufferPool.Put(buffer)

    buffer.WriteString("temporary data")

    // Process the buffer contents.
}
\end{lstlisting}

A pool can reduce allocation pressure when temporary objects are repeatedly created and discarded. It is not a permanent cache: the runtime may remove pooled objects at any time, including during garbage collection.\\

Objects obtained from a pool must be used carefully. They should be reset as necessary, should not be returned to the pool while another operation still uses them, and should not be assumed to preserve their identity across calls.\\

Pooling is most useful when profiling demonstrates that repeated allocation of suitable temporary objects is a meaningful cost.\\

\subsection{Measuring Memory Usage and Allocation Behaviour}

Go includes several tools for observing memory allocation and garbage collection. These facilities are essential when diagnosing unexpectedly high memory consumption or poor performance.\\

\subsubsection{Runtime Memory Statistics}

The \texttt{runtime.ReadMemStats} function provides information about heap allocation, heap objects and garbage-collection activity.\\

\begin{lstlisting}
package main

import (
    "fmt"
    "runtime"
)

func main() {
    var stats runtime.MemStats

    runtime.ReadMemStats(&stats)

    fmt.Printf(
        "Heap allocated: %d bytes\n",
        stats.HeapAlloc,
    )

    fmt.Printf(
        "Total allocated: %d bytes\n",
        stats.TotalAlloc,
    )

    fmt.Printf(
        "Heap objects: %d\n",
        stats.HeapObjects,
    )

    fmt.Printf(
        "GC cycles: %d\n",
        stats.NumGC,
    )
}
\end{lstlisting}

The statistics have different meanings:

\begin{itemize}
    \item \texttt{HeapAlloc} reports bytes of allocated heap objects at the time of the snapshot.
    \item \texttt{TotalAlloc} reports cumulative bytes allocated for heap objects since program start, including allocations for objects that have subsequently been collected.
    \item \texttt{HeapObjects} reports the number of allocated heap objects.
    \item \texttt{NumGC} reports the number of completed garbage-collection cycles.
\end{itemize}

These values should be interpreted together. A large cumulative allocation count does not necessarily indicate a memory leak, since a program may allocate and discard temporary objects frequently.\\

Similarly, the amount of live heap memory does not equal the total memory used by the process. Runtime metadata, goroutine stacks, executable mappings and other memory allocations may contribute to the process's overall footprint.

\subsubsection{Benchmarking Allocations}

The \texttt{testing} package can report allocation metrics during benchmarks.

\begin{lstlisting}
func BenchmarkProcess(b *testing.B) {
    b.ReportAllocs()

    for i := 0; i < b.N; i++ {
        processData()
    }
}
\end{lstlisting}

The benchmark can be run using:

\begin{lstlisting}
go test -bench=. -benchmem
\end{lstlisting}

The \texttt{-benchmem} option reports metrics such as bytes allocated per operation and allocations per operation. These measurements can help determine whether a proposed optimisation reduces allocation overhead.\\

Benchmark results should be collected under representative conditions. Compiler optimisations, CPU architecture, Go version, input size and concurrent workload can all affect the results.\\

\subsubsection{Heap Profiling}

Go's \texttt{runtime/pprof} and \texttt{net/http/pprof} packages provide profiling support. Heap profiles can identify functions responsible for allocating memory and help distinguish high allocation rates from long-lived retained objects.\\

For example, a program can write a heap profile to a file:

\begin{lstlisting}
package main

import (
    "os"
    "runtime/pprof"
)

func writeHeapProfile() error {
    file, err := os.Create("heap.prof")
    if err != nil {
        return err
    }
    defer file.Close()

    return pprof.WriteHeapProfile(file)
}
\end{lstlisting}

The resulting profile can be examined using the Go profiling tools, for example:

\begin{lstlisting}
go tool pprof heap.prof
\end{lstlisting}

Profiling information can be used to investigate allocation sites, retained memory, and object lifetimes. In particular, profiles help determine whether high memory usage results from temporary allocations, long-lived caches, large backing arrays, or other retained references.\\

Profiling should precede substantial optimisation work because it provides evidence about the actual behaviour of the application rather than relying on assumptions about which code is expensive.\\

\subsection{Finalizers and Explicit Resource Management}

Go provides runtime finalisers through facilities in the \texttt{runtime} package, including \texttt{runtime.SetFinalizer}. A finaliser can be associated with an object and may be executed after the object becomes unreachable.\\

Finalisers are generally unsuitable as the primary mechanism for resource management. Their execution is not guaranteed to occur promptly, and program termination may occur before a finaliser has run. Finalisers also introduce complications involving object reachability, ordering, and concurrency.\\

For ordinary resources such as files, network connections, and locks, explicit lifecycle management is preferable.

For example:

\begin{lstlisting}
func readFile(name string) error {
    file, err := os.Open(name)
    if err != nil {
        return err
    }
    defer file.Close()

    // Perform file operations here.
    return nil
}
\end{lstlisting}

The use of \texttt{defer} makes the intended cleanup operation clear and ensures that it runs when the function returns, including when the function returns through an error path.\\

Memory management and resource management are related but distinct concerns. Garbage collection manages the lifetime of unreachable Go heap objects, while application code must still manage resources whose correct release depends on external systems or explicit lifecycle rules.

\subsection{Conclusion}

Go provides automatic memory management through a combination of compiler analysis, runtime allocation mechanisms, and a concurrent tracing garbage collector. Escape analysis determines whether values can remain in stack storage or require heap allocation, while the garbage collector identifies unreachable heap objects and reclaims their storage.\\

The language's memory model supports pointers, slices, maps, interfaces, and dynamically sized data structures without requiring ordinary application code to perform explicit memory deallocation. At the same time, these abstractions have important consequences for allocation behaviour, object reachability, and memory retention.\\

Efficient memory management in Go therefore depends on understanding more than the existence of a garbage collector. Developers should consider allocation rates, data-structure design, goroutine lifetimes, backing-array retention, external resource management, and the behaviour of the runtime under realistic workloads.\\

The standard library provides useful mechanisms for observing and controlling memory behaviour, including runtime statistics, garbage-collection settings, benchmarking, and profiling. By combining these facilities with clear object-lifetime policies and measurement-driven optimisation, developers can build Go applications that use memory efficiently while retaining the safety and simplicity of automatic memory management.\\